\let\negmedspace\undefined
\let\negthickspace\undefined
\documentclass[journal]{IEEEtran}
\usepackage[a5paper, margin=10mm, onecolumn]{geometry}
%\usepackage{lmodern} % Ensure lmodern is loaded for pdflatex
\usepackage{tfrupee} % Include tfrupee package

\setlength{\headheight}{1cm} % Set the height of the header box
\setlength{\headsep}{0mm}     % Set the distance between the header box and the top of the text

\usepackage{gvv-book}
\usepackage{gvv}
\usepackage{cite}
\usepackage{amsmath,amssymb,amsfonts,amsthm}
\usepackage{algorithmic}
\usepackage{graphicx}
\usepackage{textcomp}
\usepackage{xcolor}
\usepackage{txfonts}
\usepackage{listings}
\usepackage{enumitem}
\usepackage{mathtools}
\usepackage{gensymb}
\usepackage{comment}
\usepackage[breaklinks=true]{hyperref}
\usepackage{tkz-euclide} 
\usepackage{listings}
% \usepackage{gvv}                                        
\def\inputGnumericTable{}                                 
\usepackage[latin1]{inputenc}                                
\usepackage{color}                                            
\usepackage{array}                                            
\usepackage{longtable}                                       
\usepackage{calc}                                             
\usepackage{multirow}                                         
\usepackage{hhline}                                           
\usepackage{ifthen}                                           
\usepackage{lscape}
\begin{document}

\bibliographystyle{IEEEtran}

\title{5.9.8}
\author{EE25BTECH11054 - S.Harsha Vardhan}
% \maketitle
% \newpage
% \bigskip
\maketitle 
\renewcommand{\thefigure}{\theenumi}
\renewcommand{\thetable}{\theenumi}
\setlength{\intextsep}{10pt} % Space between text and floats

\numberwithin{align}{enumi}
\numberwithin{figure}{enumi}
\renewcommand{\thetable}{\theenumi}
\vspace{-1em}
\textbf{Question 5.9.8}
A father's age is three times the sum of the ages of his two children. After 5 years his age will be two times the sum of their ages. Find the present age of the father.

\solution
Let the present age of the father be $F$.
Let the sum of the present ages of his two children be $S$.
\begin{align}
F &= 3S  \\
F - 3S &= 0 \\
F + 5 &= 2(S + 10)  \\
F + 5 &= 2S + 20  \\
F - 2S &= 15 
\end{align}

We now have a system of two linear equations with two variables, $F$ and $S$. In matrix form $\mtx{A}\vec{x} = \vec{b}$, this is:
\begin{align}
\myvec{1 & -3 \\ 1 & -2} \myvec{F \\ S} &= \myvec{0 \\ 15} 
\end{align}
To solve this system, we use Gaussian elimination on the augmented matrix $\augvec{1}{1}{\vec{A}&\vec{b}}$.
\begin{align}
\augvec{2}{1}{1 & -3 & 0 \\ 1 & -2 & 15} \xrightarrow{R_2 \to R_2 - R_1}
\augvec{2}{1}{1 & -3 & 0 \\ 0 & 1 & 15} \xrightarrow{R_1 \to R_1 + 3R_2}\augvec{2}{1}{1 & 0 & 45 \\ 0 & 1 & 15} 
\end{align}

\begin{align}
S = 15 \tag{9}
\end{align}
 
From the reduced row echelon form, we get the solution:
\begin{align}
F = 45 \quad \text{and} \quad S = 15 \tag{11}
\end{align}
The question asks for the present age of the father.
Thus, the present age of the father is 45 years.
\begin{figure}[h]
    \centering
    \includegraphics[width=1.1\linewidth]{figs/Figure_7.png}
    \caption{}
    \label{fig:placeholder}
\end{figure}

\end{document}